\chapter*{Introduction}

Antilinear representations play an important role in algebra, physics and topology \cite{AtiSeg, AtiKtheory, NeebOlaf, DYSON, BradDav, FreMoo, TopIns, Bootstrap}. Besides physics, they also appear in chemistry and crystallography \cite{Bro, Ceulemans, JensenSaue}. They were first introduced by Wigner \cite{WIG} under the name of co-representations to model time-reversal symmetry in quantum mechanics. 

Algebraically, for $G$ a group with an index two subgroup $H$, an antilinear representation of $G$ is an action of $G$ on a complex vector space such elements of $H$ act linearly, whilst elements of $G\setminus H$ act antilinearly. The key result in the theory is {\em Dyson's Tenfold Way}, which classifies irreducible antilinear representations into 10 fundamentally different classes. The goal of this text is to generalise Dyson's Tenfold Way to the following setting:
\begin{itemize}
\item the group $G$ is any (discrete) group, 
\item instead of complex vector spaces, we work with any pair $(\bL, \tau)$, where $\bL$ is a field and $\tau$ is an order two automorphism of $\bL$. We study actions of $G$ on finite-dimensional vector spaces over $\bL$, such that elements of $H$ act linearly and elements of $G \setminus H$ act $\tau$-linearly. We write
  \[ \Irr_{\bbL}^{\hspace{0.1em} \star}(G)\]
  for the set of (isomorphism classes of) such irreducible representations of $G$ over $\bL$.
\end{itemize}

The main outcome of our work is that we show these representations fall into different {\em types}, that capture key algebraic properties of a representation. For $\bL = \bC$, we recover Dyson's Tenfold Way. 

As we will use the term ``type'' in a few different ways, let us summarise, without giving precise definitions, the various notations of a type we will use. For $V$ a representation, we have:
\begin{itemize}
\item {\em The type} $\fT(V)$. This is the main overriding type that keeps track of all the algebraic information about $V$.
As a result of our classification, there are 24 types; the reader can find them in Table~\ref{tab:24types}.
By $\fT(V)$ we understand either the number from 1 to 24, or the actual information in the corresponding row, some sort of a ``barcode''.
\item {\em The D-type} $\fD(V)$, or {\em Dyson type}. 
This is an alternative name for the type $\fT(V)$ which we use in certain situations, such as when $\bL$ is algebraically closed, when only 12 types are possible.
$\fD(V)$ is the Roman numeral in the leftmost column of Table~\ref{tab:24types}.
It is called the Dyson type because when $\bL = \bC$, only 10 of the 12 are possible, and $\fD(V)$ is simply the label of Dyson's Tenfold Way.
\item {\em The S-type} $\fS(V)$, or {\em semilinear type}.
This is the information in the column $\bL[H] \subset \bL_{\star \hspace{-0.1em}}[G]$ of Table~\ref{tab:24types}.
There are 3 possible S-types: \textbf{A}, \textbf{B} or \textbf{C}.
\item {\em The C-type} $\fC(V)$, or {\em classical type}.
This is the information in the column $\bK[H] \subset \bL[H]$ of Table~\ref{tab:24types}.
There are 3 possible C-types: \textbf{A}, \textbf{B} or \textbf{C}.
\item {\em The L-type} $\fL(V)$, or {\em linear type}.
  This is the information in the four columns
of Table~\ref{tab:24types}
that do not contain  $\bL_{\star \hspace{-0.1em}}[G]$. By inspection,  
there are 18 possible L-types.
\end{itemize}  
For instance, if $\fT(V)=5$, then we examine row 5 of Table \ref{tab:24types} and conclude that
\[ \fD(V) = \mbox{II}, \ \ \fS(V) = \textbf{C}, \ \ \fC(V) = \textbf{A}, \ \ \fL(V) =( \textbf{B},\textbf{A},\textbf{1},\textbf{3(b)}).\]
Now we outline the main ideas behind these notions of types.

\section*{D-Type}

Let us first recall, through our lens of types, the theory of $\bR$-linear and antilinear representations, and Dyson's Tenfold Way. The types that play a role here are the C-type, the S-type, and the D-type.

\subsection*{C-Type and $\bR$-Linear Representations}

Let $H$ be a (discrete) group.
Let us examine the relationship between the sets
$\Irr_{\bR}(H)$ and $\Irr_{\bC}(H)$
of finite-dimensional irreducible $\bR$-linear and $\bC$-linear representations.
Explicitly, for any $V \in \Irr_{\bC}(H)$, there is a unique $X \in \Irr_{\bR}(H)$ with $\Hom_{\bR[H]}(X, V|_{\bR}) \neq 0$.
For this $X$, the endomorphism ring $\End_{\bR[H]}(X)$ is a finite-dimensional division algebra over $\bR$.
By Frobenius' Theorem there are three possibilities:

\begin{itemize}
    \item[(\textbf{A})] $\End_{\bR[H]}(X) \cong \bR$, $X_{\bC} \cong V$, and $V|_{\bR} \cong X \oplus X$.
    \item[(\textbf{B})] $\End_{\bR[H]}(X) \cong \bC$, $X_{\bC} \cong V \oplus \overline{V}$, and $V|_{\bR} \cong X$.
    \item[(\textbf{C})] $\End_{\bR[H]}(X) \cong \bH$, $X_{\bC} = V \oplus V$, and $V|_{\bR} \cong X$.
\end{itemize}

In our language of types, still yet to be defined, the label (\textbf{A}), (\textbf{B}) or (\textbf{C}) that occurs is exactly the C-type $\fC(X)$ of $X$, or equally the C-type $\fC(V)$ of $V$. 

Observe that there are two different algebraic ways to determine the C-type. The first is via the endomorphism algebra $\End_{\bR[H]}(X)$. The second is through the interplay between induction and restriction functors. This interplay is the main perspective behind our theory of types.

\subsection*{The S-Type and Antilinear Representations}
Now let $G$ be a group which acts on $\bC$ non-trivially, in the sense that there is a surjective group homomorphism $\sigma \colon G \twoheadrightarrow \Gal(\bC / \bR)$.

Similarly to $\bR$-linear representations, we can relate irreducible antilinear representations $\Irr_{\bC}^{\sigma}(G)$ of $G$ to $\Irr_{\bC}(H)$, where now $H$ denotes the index two normal subgroup $H \coloneqq \ker(\sigma)$ of $G$.
Consider the induction and restriction functors
\[
(-)|_H \colon \Rep_{\bC}^{\sigma}(G) \rightarrow \Rep_{\bC}(H), \qquad \Ind_H^G(-) \colon \Rep_{\bC}(H) \rightarrow \Rep_{\bC}^{\sigma}(G),
\]
and an action of $C_2 = \langle \epsilon \rangle$ on $\Irr_{\bC}(H)$, which is defined by fixing some $g \in G \setminus H$, and for $V \in \Irr_{\bC}(H)$ setting $\epsilon * V$ to be $\overline{V}$ as a $\bC$ vector space, with action of $h \in H$ by $h * v \coloneqq (g^{-1} hg)(v)$.

For any $V \in \Irr_{\bC}(H)$, there is a unique $W \in \Irr_{\bC}^{\sigma}(G)$ with $\Hom_{\bC_{\sigma}\hspace{-0.1em}[G]}(W, \Ind_H^G(V)) \neq 0$. For this $W$, the endomorphism algebra $\End_{\bC_{\sigma}\hspace{-0.1em}[G]}(W)$ is a finite-dimensional division algebra over $\bR$, and Frobenius' Theorem again results in three possibilities:
\begin{itemize}
    \item[(\textbf{A})] $\End_{\bC_{\sigma}\hspace{-0.1em}[G]}(W) \cong \bR$, $W|_H \cong V$, and $\Ind_H^G(V) \cong W \oplus W$.
    \item[(\textbf{B})] $\End_{\bC_{\sigma}\hspace{-0.1em}[G]}(W) \cong \bC$, $W|_H \cong V \oplus \epsilon * V$, and $\Ind_H^G(V) \cong W$.
    \item[(\textbf{C})] $\End_{\bC_{\sigma}\hspace{-0.1em}[G]}(W) \cong \bH$, $W|_H \cong V \oplus V$, and $\Ind_H^G(V) \cong W$.
\end{itemize}
In our language, this label is exactly the S-type $\fS(W)$ of $W$, or equally the S-type $\fS(V)$ of $V$.

\subsection*{The D-Type and Dyson's Tenfold Way}
Given a pair $(G,\sigma)$ as before, let us consider both interplays above simultaneously
\[
\Irr^{\sigma}_{\bC}(G) \longleftrightarrow \Irr_{\bC}(H)\longleftrightarrow \Irr_{\bR}(H), \quad 
W \leftrightarrow V \leftrightarrow X.
\]
In this situation, it is natural to consider how $W$ and $X$ relate to one another. For example, to ask how the S-type for the pair $(V,W)$ is related to the C-type for the pair $(V,X)$.

To answer this, another linear representation category $\Rep_{\bR}(G)$ enters the scene.
There are induction and restriction functors now between the four categories: 
% https://q.uiver.app/#q=WzAsNCxbMSwwLCJcXFJlcF97XFxiUn0oRykiXSxbMiwxLCJcXFJlcF97XFxiQ31ee1xcc2lnbWF9KEcpIl0sWzEsMiwiXFxSZXBfe1xcYkN9KEgpIl0sWzAsMSwiXFxSZXBfe1xcYlJ9KEgpIl0sWzAsMV0sWzEsMF0sWzIsMV0sWzEsMl0sWzAsM10sWzMsMl0sWzIsM10sWzMsMF1d
\begin{equation} \label{2categories}
  \begin{tikzcd}[sep=2.25em]
	& {U \in \Rep_{\bR}(G)} & \\
	{X\in \Rep_{\bR}(H)} && {\Rep_{\bC}^{\sigma}(G) \ni W} \\
	& {V\in \Rep_{\bC}(H)}
	\arrow[from=1-2, to=2-1]
	\arrow[from=1-2, to=2-3]
	\arrow[from=2-1, to=1-2]
	\arrow[from=2-1, to=3-2]
	\arrow[from=2-3, to=1-2]
	\arrow[from=2-3, to=3-2]
	\arrow[from=3-2, to=2-1]
	\arrow[from=3-2, to=2-3]
\end{tikzcd}\end{equation}

Let $U \in \Irr_{\bR}(G)$ such that $\Hom_{\bR[G]}(U, W|_{\bR}) \neq 0$, or equivalently that $\Hom_{\bR[G]}(U, \Ind_{H}^G(X)) \neq 0$. This $U$ is unique up to replacement with $U \otimes \pi$, the tensor product of $U$ with the sign representation $\pi$ of $G$ over $\bR$ (with kernel $H$). For either choice, we have that $\End_{\bR[G]}(U) = \End_{\bR[G]}(U \otimes \pi)$. 

Now pick any $V \in \Irr_{\bC}(H)$, then choose $X, W$ and $U$ as above. We mark where these modules belong on the diagram~\eqref{2categories} for the convenience of the reader. The S-type and the C-type of V can be recovered from the the triple 
\[
(\End_{\bR[H]}(X), \End_{\bR[G]}(U), \End_{\bC_{\sigma}\hspace{-0.1em}[G]}(W)) \in \{\bR, \bC, \bH\}^3 \, ,
\]
which is independent of the choices of $X$, $U$ and $W$. The D-type $\fD(V)$ of $V$, still yet to be defined, is an invariant which records how induction and restriction works between $V$, $X$, $W$ and $U$ in diagram~\eqref{2categories}. It is sufficiently rich to determine the endomorphism ring triple above.

The fundamental result of Dyson, phrased in our language, is the following.
\begin{theorem*}[Dyson's Tenfold Way \cite{DYSON}]
Suppose $G$ is a (discrete) group.
Given $V \in \Irr_{\bC}(H)$, there are ten possibilities for $\fD(V)$, listed with corresponding endomorphism algebra triples:
\[
\renewcommand{\arraystretch}{1.1}
\begin{array}{@{}c@{\;}l@{\qquad}
                c@{\;}l@{\qquad}
                c@{\;}l@{\qquad}
                c@{\;}l@{\qquad}
                c@{\;}l@{}}
\mathrm{I}    & (\bR,\bR,\bR) &
\mathrm{III}  & (\bR,\bR,\bC) &
\mathrm{V}    & (\bC,\bR,\bR) &
\mathrm{VII}  & (\bC,\bC,\bC) &
\mathrm{IX}   & (\bH,\bC,\bR) \\

\mathrm{II}   & (\bR,\bC,\bH) &
\mathrm{IV}   & (\bC,\bC,\bC) &
\mathrm{VI}   & (\bC,\bH,\bH) &
\mathrm{VIII} & (\bH,\bH,\bH) &
\mathrm{X}    & (\bH,\bH,\bC)
\end{array}
\]
\end{theorem*}

We thus see that the endomorphism ring triple is almost enough to determine the D-type, but not quite. The unique repeat is in D-types IV and VII. In Dyson's original notation these are called $\sC \sC 2$ and $\sC \sC 1$ respectively - he uses the pair $(\End_{\bC_{\sigma}\hspace{-0.1em}[G]}(W), \End_{\bR[H]}(X))$ as his label for the type, so II is called $\sQ \sR$ and so on (with letters denoting $\sR$eal, $\sC$omplex and $\sQ$uaternionic).

Conceptually, the main reason that there are 10 possible D-types is that the Brauer-Wall monoid $\BrW(\bR) \sqcup \BrW(\bC)$ has order 10 (see \S \ref{sect:recoverDyson10}). For us to obtain finiteness in our classification of types over general fields, it turns out that the induction-restriction perspective is more profitable, as for general fields the corresponding Brauer-Wall monoids will typically be infinite.

\section*{Semilinear Representations for Degree Two Extensions}
Let us move on from the complex numbers to more general fields.
Suppose now that $G$ is a group, and $\bL$ is a field, upon which $G$ acts through a group homomorphism $\sigma \colon G \rightarrow \Aut(\bL)$.
In this context, a \emph{semilinear representation} of $G$ over $\bL$ is a finite-dimensional $\bL$-vector space $W$, with a homomorphism $\rho \colon G \rightarrow \Aut_{\bK}(W)$, such that each $\rho(g)$ is $\sigma(g)$-linear.

If $\bL / \bK$ is finite, where we write $\bK \coloneqq \bL^G$, then automatically $\bL / \bK$ is Galois, and the map $\sigma \colon G \rightarrow \Gal(\bL / \bK)$ is surjective. Under this assumption, the category of such semilinear representations $\Rep_{\bbL}^{\sigma}(G)$ can be described in terms of $\Rep_{\bbL}(H)$, where $H \coloneqq \ker(\sigma \colon G \rightarrow \Gal(\bL/\bK))$ \cite{TAYSchur}.

\subsection*{S-Type for Degree Two Extensions} 
The focus of the present work is on understanding $\Rep_{\bbL}^{\sigma}(G)$ in the smallest non-trivial case, when $\bL / \bK$ has degree two. From now until the end of introduction, this is an assumption.
Of course, this includes the case when $\bL / \bK$ is $\bC / \bR$.

As for antilinear representations over $\bC$, there is an analogously defined action of $C_2 = \langle \epsilon \rangle$ on $\Irr_{\bL}(H)$, and semilinear induction and restriction functors
\[
(-)|_H \colon \Rep_{\bbL}^{\sigma}(G) \rightarrow \Rep_{\bbL}(H), \qquad \Ind_H^G(-) \colon \Rep_{\bbL}(H) \rightarrow \Rep_{\bbL}^{\sigma}(G).
\]
For any $V \in \Irr_{\bC}^{\sigma}(G)$, there is still a unique $W \in \Irr_{\bbL}^{\sigma}(G)$ with $\Hom_{\bbL_{\sigma}\hspace{-0.1em}[G]}(W, \Ind_H^G(V)) \neq 0$.
The crucial difference is that instead of $\End_{\bbL_{\sigma}\hspace{-0.1em}[G]}(W)$ the key quantity is the ratio:
\[
d(V,W) \coloneqq \dim_{\bbK}\End_{\bbL_{\sigma}\hspace{-0.1em}[G]}(W) / \dim_{\bbL}\End_{\bbL[H]}(V).
\]
With this notation, there are three possibilities, which we call the S-type $\fS(V) = \fS(W)$:
\begin{itemize}
    \item[(\textbf{A})] $d(V,W) = 1$, $W|_H \cong V$, and $\Ind_H^G(V) \cong W \oplus W$.
    \item[(\textbf{B})] $d(V,W) = 2$, $W|_H \cong V \oplus \epsilon * V$, and $\Ind_H^G(V) \cong W$.
    \item[(\textbf{C})] $d(V,W) = 4$, $W|_H \cong V \oplus V$, and $\Ind_H^G(V) \cong W$.
\end{itemize}

\subsection*{Classification of Semilinear Representations}
The reason for our interest in the classification into three types \textbf{A}, \textbf{B} and \textbf{C}, is that it provides a method to construct all irreducible semilinear representations of $G$ over $\bL$ from $\Irr_{\bbL}(H)$, provided that one already knows the S-type of each $V \in \Irr_{\bbL}(H)$.
Concretely, pick $V \in \Irr_{\bbL}(H)$.

\vspace{0.2em}

\noindent \hspace{0.7em} $\bullet$ If $\fS(V)=$ \textbf{A}, then $V$ extends uniquely to an irreducible semilinear representation of $G$ over $\bL$.

\noindent \hspace{0.7em} $\bullet$ If $\fS(V)\in\{\textbf{B},\textbf{C}\}$, then $\Ind_H^G(V)$ is an irreducible semilinear representation of $G$ over $\bL$.

\vspace{0.2em}

\noindent Furthermore, all $W \in \Irr_{\bbL}^{\sigma}(G)$ arise this way.

The question is thus how to determine, given $V \in \Irr_{\bbL}(H)$, whether $V$ has type \textbf{A}, \textbf{B} or \textbf{C}.
We can immediately see that the $C_2$-action allows us to distinguish S-type \textbf{B}:
\[
\text{\centering{$\fS(V)=\textbf{B}$ if and only if $V \not\cong \epsilon * V$.}}
\]
The following question is therefore of fundamental interest:

\begin{question1*}
For $V \in \Irr_{\bbL}(H)$ with $V \cong \epsilon * V$, how can one tell if $V$ has S-type \emph{\textbf{A}} or \emph{\textbf{C}}?
\end{question1*}

\section*{The Type and L-Type of $V$}
How is Question 1 answered when $\bL / \bK$ is $\bC / \bR$?
This is where Dyson's Tenfold Way enters the picture.
If $V \in \Irr_{\bbL}(H)$ has $V \cong \epsilon * V$, then $\End_{\bC_{\sigma}\hspace{-0.1em}[G]}(W)$ is either $\bR$ or $\bH$. Examining the possibilities in Dyson's Tenfold Way above, we see that 
\[
\fS(V) = \twopartdef{\text{\textbf{A}}}{(\End_{\bR[H]}(X), \End_{\bR[G]}(U)) \in \{(\bR, \bR), (\bC, \bR), (\bH, \bC)\}}{\text{\textbf{C}}}{(\End_{\bR[H]}(X), \End_{\bR[G]}(U)) \in \{(\bR, \bC), (\bC, \bH), (\bH, \bH)\}}.
\]
Therefore Question 1 is answered by examining the pair $(\End_{\bR[H]}(X), \End_{\bR[G]}(U))$, a purely linear representation-theoretic quantity, which is computable using classical (linear) methods.

For general fields, we want to understand if something similar is possible in general. Is there some linear data we can associate to $V$, which determines the S-type of $V$?

Besides the four categories on the diagram~\eqref{2categories},
we can bring one more category into play, arriving at {\em the diamond}
\begin{equation}\label{eq:maincats}
\begin{tikzcd}
	& {\Rep_{\bbK}(G)} \\
	{\Rep_{\bbL}(G)} & {\Rep_{\bbK}(H)} & {\Rep_{\bbL}^{\hspace{0.1em}\sigma}(G)} \\
	& {\Rep_{\bbL}(H)}
	\arrow["3"', tail reversed, from=1-2, to=2-1]
	\arrow["4", tail reversed, from=1-2, to=2-3]
	\arrow["4", tail reversed, from=2-2, to=1-2]
	\arrow["3"', tail reversed, from=2-2, to=3-2]
	\arrow["4", tail reversed, from=3-2, to=2-1]
	\arrow["3"', tail reversed, from=3-2, to=2-3]
\end{tikzcd}
\end{equation}
of categories, connected by canonical induction and restriction functors. 

Let us from now on assume that $\charfield(\bK) \neq 2$.
We already know that there are three possibilities for induction and restriction between $\Rep_{\bbL}(H)$ and $\Rep_{\bbL}^{\hspace{0.1em}\sigma}(G)$, recorded in the S-type. It turns out that this finiteness extends to the other pairs: there are either three or four possibilities for induction and restriction between each category in adjacent pair of categories in diagram~\eqref{eq:maincats}.

We label each extension with the exact number. When there are three possibilities, these are called \textbf{A}, \textbf{B} and \textbf{C}, and when there are four these are called $\mathbf{1}, \mathbf{2}$, \textbf{3(a)} and \textbf{3(b)}.

\begin{definition*}
The \emph{type} of $V$ is the tuple
\[
\fT(V) \in \{\mathbf{A},\mathbf{B},\mathbf{C}\}^3 \times \{\text{\textbf{1}}, \text{\textbf{2}}, \text{\textbf{3(a)}}, \text{\textbf{3(b)}}\}^3
\]
which records induction and restriction in diagram~\eqref{eq:maincats}, read from left to right.
\end{definition*}

The S-type is the third entry of $\fT(V)$, recording induction and restriction between $\Rep_{\bbL}(H)$ and $\Rep_{\bbL}^{\hspace{0.1em}\sigma}(G)$. The C-type is the second entry of $\fT(V)$, recording induction and restriction between $\Rep_{\bbL}(H)$ and $\Rep_{\bbK}(H)$. 

\begin{definition*}
The \emph{linear type}, or L-type, of $V$ is the tuple
\[
\fL(V) \in \{\mathbf{A},\mathbf{B},\mathbf{C}\}^2 \times \{\text{\textbf{1}}, \text{\textbf{2}}, \text{\textbf{3(a)}}, \text{\textbf{3(b)}}\}^2
\]
obtained from $\fT(V)$ by forgetting the entries which involve $\Rep_{\bbL}^{\hspace{0.1em}\sigma}(G)$.
\end{definition*}

By definition, the L-type consists of exactly the part of $\fT(V)$ which involves only classical (linear) representation theory.
Using $\fT(V)$ and $\fL(V)$, we can make Question 1 more precise.

\begin{question2*}
    For $V \in \Irr_{\bbL}(H)$, can $\fT(V)$ be recovered from $\fL(V)$?
\end{question2*}

\section*{Classification of Types}

Our first main result is a generalisation of Dyson's Tenfold Way to arbitrary degree two extensions, which yields a complete answer to Question 2.

\begin{theoremA}
For $V \in \Irr_{\bL}(H)$, there are 24 possibilities for $\fT(V)$, listed in Table \ref{tab:24types}.
\end{theoremA}

Examining the 24 types listed in Table \ref{tab:24types}, we see that there are 18 possible values for $\fL(V)$, and therefore that $\fL(V)$ does not determine $\fT(V)$ in general. This resolves Question~2.

In many cases, however, Theorem A does allow us to deduce $\fT(V)$ from $\fL(V)$. Examining the 24 semilinear types listed in Table \ref{tab:24types} in more detail, we see that
\begin{itemize}
    \item in 12 out of the 24 types, the L-type $\fL(V)$ \emph{does} determine $\fT(V)$, however, and
    \item in the remaining 12 types, the L-type $\fL(V)$ determines two possibilities for $\fT(V)$.
\end{itemize}

\section*{Applications of the Classification of Types}
The significance of Theorem A is that it provides a means by which we can prove results for semilinear representations using results about linear representations. Let us now describe these, together with other applications of Theorem A.

\subsection*{Symmetric Groups}
The possible types $\fT(V)$
are often significantly restricted when one considers specific groups.
For example, suppose that $\bL / \bK$ has characteristic $0$
and the group-subgroup pair is $A_n \leq S_n$, an alternating/symmetric group with $n \geq 3$.

Given $V \in \Irr_{\bL}(A_n)$, using Theorem A one can show that the possible types $\fT(V)$ are
only 1, 2, 3 and 12 of Tables~\ref{tab:24types5} and \ref{tab:24types6} (Corollary \ref{cor:symgroups}).
In particular, the L-type $\fL(V)$ determines $\fT(V)$.

More importantly,  the only possible S-types $\fS(V)$ are \textbf{A} or \textbf{B}.
Hence, the semilinear representation theory of $S_n$ over $\bL$ is completely determined by $\Irr_{\bL}(A_n)$ and its action by $C_2$.
In particular, everything is completely determined by ordinary character theory. For $V \in \Irr_{\bL}(A_n)$:

\vspace{0.2em}

\noindent \hspace{0.7em} $\bullet$ If $V \cong \epsilon * V$, then $V$ extends uniquely to an irreducible semilinear representation of $S_n$ over $\bL$.

\noindent \hspace{0.7em} $\bullet$ If $V \not\cong \epsilon * V$, then $\Irr_H^G(V)$ is an irreducible semilinear representation of $S_n$ over $\bL$.

\noindent \hspace{0.7em} $\bullet$ All irreducible semilinear representations of $S_n$ over $\bL$ arise this way.

\subsection*{Cyclotomic Fields}
The possible types $\fT(V)$ are also often significantly restricted when one concentrates on a specific extension $\bL / \bK$.
For example, if $\bL$ is a finite field, then its Brauer group is trivial.
It follows that $\fC(V)$ and $\fS(V)$ are never \textbf{C} \cite[Cor.\ 9.4]{TAYSchur}.
Hence, by Theorem~A only 10 types $\fT(V)$ are possible, as seen by consulting Table \ref{tab:24types}.

Another important extension is $\bC / \bR$.
Here, by Dyson's Tenfold Way, only 10 types are possible. In Section~\ref{sect:recoverDyson10} we give a short proof of Dyson's Tenfold Way using Theorem~A, and show how in light of Theorem~A it follows directly from the classification of $C_2$-graded division algebras over $\bR$.

Another interesting class of extensions are the cyclotomic extensions.
Associated to the pair $(G, \sigma)$, we can consider the extension $\bQ(\zeta_e)_+ \leq \bQ(\zeta_e)$, where $e$ denotes the exponent of $H = \ker(\sigma)$ (assumed large enough so that this is a degree two extension). Our second main result is the following.

\begin{theoremB}
There are $15$ possibilities for the type $\fT(V)$ over $\bQ(\zeta_e)_+ \leq \bQ(\zeta_e)$, listed in Table~\ref{tab:15typesQe}.
\end{theoremB}

Examining in more detail the possible types, this allows us to construct all irreducible semilinear representations of $G$ over $\bQ(\zeta_e)_+ \leq \bQ(\zeta_e)$ using only the complex character table of $H$ and its action by $C_2$, together with the complex character table of $G$ and its action by $\Gal(\bQ(\zeta_{2e}) / \bQ(\zeta_e))$.

\section*{Other Results}
\subsection*{Antilinear Induction Theorem}

The crux of the proof of Theorem B relies on the following induction theorem for antilinear representations (semilinear representations for the extension $\bC / \bR$).

\begin{theoremC}
Suppose that $G$ is finite and $W \in \Irr_{\bC}^{\sigma}(G)$. Then there are subgroups $H_i \leq H$ and $G_j \leq G$ with each $G_j \not\subset H$, one-dimensional representations $V_i \in \Irr_{\bC}(H_i)$ and $W_j \in \Irr_{\bC}^{\sigma}(G_i)$, and $m_i, n_j \in \bZ$, such that
\[
    [V] = \sum_i m_i \cdot [\Ind_{H_i}^G(V_j)] + \sum_j n_j \cdot [\Ind_{G_j}^G(W_j)]
\]
in $K_0(\Rep_{\bC}^{\sigma}(G))$.
\end{theoremC}

We note that Theorem C requires considering induction from \emph{all} subgroups of $G$, not just the non-trivially graded subgroups $G_j$, and therefore blends the $\bC$-linear and antilinear theories.

Along with Brauer's induction Theorem, which concerns $\bC$-linear representations, Serre proved an induction theorem for $\bR$-linear representations, which states that any irreducible $\bR$-linear representation of a finite group is induced from representations of subgroups, which themselves are inflated from irreducible representations of cyclic or dihedral groups \cite[Lem.\ 73.19]{CR2}.

Theorem C follows from a slightly stronger statement we prove, Theorem \ref{thm:antilinearinductionthm}, which also applies to groups which also act trivially on $\bC$, and can be viewed as a simultaneous generalisation of both Brauer's Induction Theorem and Serre's Induction Theorem (see Section \ref{sect:recoveryofclassical}).

This generalisation furthermore explains the appearance of dihedral groups in Serre's Induction Theorem. Indeed, viewing each dihedral group $D_{2n}$ with it's natural action on $\bC$ with kernel $C_n$, dihedral groups are characterised as exactly those finite groups $G$ which admit a faithful $1$-dimensional antilinear representation (Lemma \ref{lem:1dimfftypeRisdihedral}).

From Theorem C, we also obtain the consequence that all antilinear representations admit models over $\bQ(\zeta_e)$ (Theorem \ref{thm:antilinearcyclotomicfinal}). Similarly, this can be viewed as a generalisation of the fact that any $V \in \Irr_{\bC}(H)$ admits a model over $\bQ(\zeta_e)$ and that any $X \in \Irr_{\bR}(H)$ admits a model over $\bQ(\zeta_e)_+$ \cite{PAS}.

\subsection*{Classification of Real Structures}

For a fixed extension $\bL / \bK$, to give a surjective homomorphism $\sigma \colon G \rightarrow \Gal(\bL / \bK)$ is the same as to give an subgroup $H \leq G$ of index two.

One can instead shift perspective from $G$ and begin with a group $H$, and consider \emph{Real structures} on $H$, which are pairs $(G,\iota)$ where $G$ is a group and $\iota \colon H \hookrightarrow G$ is an injective group homomorphism with $\iota(H)$ index two in $G$. For example, in this way one can view the $\bR$-linear representation theory of $H$ as just one among the possible antilinear representation theories $\Rep_{\bC}^{\sigma}(G)$.
One simply considers $G = H \times C_2$,  the trivial Real structure.

In Chapter \ref{chapter:realstructures}, we classify Real structures on a group $H$, first in a explicit and elementary fashion, and them cohomologically. We relate these approaches and provide algorithms for enumerating all Real structures on a given group $H$. 


\subsection*{Classification of $C_2$-Modules}
The cohomological approach to the classification of Real structures involves the cohomology for the group $C_2$.
With this in mind, we prove a classification theorem for finitely-generated $\bZ[C_2]$-modules. We show that this category satisfies the Krull-Schmidt Theorem, meaning every object is a direct sum of indecomposable objects. We then describe the indecomposable objects in terms of certain finite linear data over the field $\bF_2$, which is algorithmically computable.

\subsection*{Antilinear Representations of Magnetic Groups}
Antilinear representations are of interest in Crystallography \cite{Bro, Wills}.
In Chapter~\ref{chap:physics} we consider \emph{magnetic groups}, which are Real structures on crystallographic groups.
We consider the 2-dimensional case, of Real structures on wallpaper groups.

We rederive the classification of 80 $2$-dimensional magnetic groups from our classification of Real structures and provide an in-depth analysis of their antilinear representation theory. In particular, we extend the theory of the Frobenius-Schur indicator to this setting, show that Dyson's Tenfold Way applies in this setting, and study the types of their irreducible representations. We calculate all relevant invariants of these groups and provide full algorithms in magma to justify the correctness of the tabulated results.

\section*{Loose Ends}
There are many aspects of our theory of types that deserve further explanation and additional exploration.
We touch upon some of them but not in sufficient depth.

\subsection*{Brauer-Wall Groups}
While we discuss them briefly in Appendix~\ref{appendix:gradedalgebras}, they deserve much more thorough investigation.
Consider a field $\bK$ of characteristic zero, its finite extension $\bF\geq \bK$ and an element of the Brauer-Wall group $A\in\BrW(\bK)$.
Can one describe all $C_2$-graded groups $G$ and quadratic extensions $\bL\geq \bK$ that ``realise'' $A$, in a sense of existence of
a surjective $C_2$-graded homomorphism $\bL_{\star \hspace{-0.1em}}[G] \twoheadrightarrow M_n (A)$?
What possible types $\fT (V)$ can the corresponding antilinear representation have? These questions appear intriguing.

\subsection*{Topological Groups}
The classical Dyson's Tenfold Way works equally well in the context of unitary representations of compact groups.
This obviously covers the case of finite groups.

Our work extends the Tenfold Way to finite-dimensional (discrete) representations of infinite groups. The next obvious step to try locally compact groups, and continuous representations.

Furthermore, we feel that the Tenfold Way is quite a general phenomenon that may work in several, more abstract contexts: more general topological groups, groups in tensor categories, topological vectors spaces, ind-vector spaces, etc.

\subsection*{Higher Dimensional Crystallographic Groups}
Chapter \ref{chap:physics} and Appendix~\ref{appendix_wallpaper} focus on the wallpaper groups, 2-dimensional crystallographic groups. Yet almost all of its provisions work in higher dimensions, and more generally for any virtually free abelian group. It would be interesting to see this theory applied in the 3-dimensional case, to classify antilinear representations of 3-dimensional magnetic groups.

\subsection*{Other Galois Groups}
Is the group $C_2$ special? What about more general Galois groups?
It feels like the next step is to consider groups $G$ acting on fields $\bL$ such that
the image of $G$ in the Galois group is cyclic of prime order. We believe that much of our theory can be done in this generality.

A much harder proposition is the question of more general Galois groups.
For instance, it would be interesting to investigate the case when the image of $G$ in the Galois group is a finite simple group.

\subsection*{Topology and Physics}
The authors are neither physicists nor topologists. Nevertheless, the theory developed here appears to have substantive consequences in both areas. We hope that our work will spark interest in specialists in those fields, and that these directions be explored further.

\section*{Acknowledgements and Author Statements}

\subsection*{Acknowledgements} The second author was supported by the departmental grant Progetto Sviluppo Dipartimentale - UNIPD PSDIP23O88 at the University of Padova, and the Heilbronn Institute of Mathematical Research at the University of Bristol.

\subsection*{AI Usage Statement} The key ideas and main results of this paper were all established without the use of LLMs, whilst free versions of LLMs available online were used for certain tasks, such as writing the associated magma code or searching the literature. The paper was entirely written without the use of LLMs.